\section{The conformal square map: illuminating but not the final solver}
\label{sec:conformal}

Treat the torque plane as the complex plane:
\begin{equation}
 \gls{sym:w}=\gls{sym:zpsi}+\mathrm j\gls{sym:zq},\qquad
 \gls{sym:zeta}=\gls{sym:w}^2.
\end{equation}
Exact expansion gives
\begin{equation}
 \operatorname{Re}\zeta=z_\psi^2-z_q^2,\qquad
 \operatorname{Im}\zeta=2z_\psi z_q=\frac{2M}{\kappa}.
 \label{eq:square-map}
\end{equation}
Constant-torque hyperbolas therefore map to horizontal lines. A loss circle
\(|w|^2=P\) maps to a circle of radius \(|\zeta|=P\). At the sample points
\(w=re^{\mathrm j\theta}\), the map gives
\(\zeta=r^2e^{\mathrm j2\theta}\): radii square and angles double.

The derivative is \(2w\), which is nonzero away from the origin. Locally, the
map is therefore a rotation and positive scaling, so it preserves the angle
between intersecting smooth curves. ``Conformal'' means local angle
preservation, not global distance or area preservation. The points \(w\) and
\(-w\) share the same square; the map is two-to-one and a current-sign branch
must be selected on inversion. Given \(\zeta=X+\mathrm jY\), a square-root
reconstruction uses
\begin{equation}
 z_\psi=\pm\sqrt{\frac{|\zeta|+X}{2}},\qquad
 z_q=\sgn(Y)\sqrt{\frac{|\zeta|-X}{2}},
 \label{eq:square-root-inverse}
\end{equation}
with the coupled sign chosen for the desired current branch. This is friendly
to fixed-function hardware: additions, one magnitude square root, and two
component square roots.

\begin{figure}[tb]
\centering
\begin{tikzpicture}
\begin{groupplot}[group style={group size=2 by 1,horizontal sep=12mm},
 axis main,axis equal image,width=.46\linewidth,height=5.6cm]
\nextgroupplot[xlabel={$z_\psi$},ylabel={$z_q$},xmin=-2.4,xmax=2.4,ymin=-2.2,ymax=2.2,title={$w$ plane}]
 \foreach \c in {.45,.9}{
  \addplot[torque curve,domain=.23:2.3,samples=80]{\c/x};
  \addplot[torque curve,domain=-2.3:-.23,samples=80]{\c/x};}
 \addplot[loss curve,domain=0:360,samples=100] ({1.55*cos(x)},{1.55*sin(x)});
\nextgroupplot[xlabel={$\operatorname{Re}\zeta$},ylabel={$\operatorname{Im}\zeta$},xmin=-2.7,xmax=2.7,ymin=-2.7,ymax=2.7,title={$\zeta=w^2$ plane}]
 \addplot[torque curve] coordinates {(-2.5,.9) (2.5,.9)};
 \addplot[torque curve,Purple!55] coordinates {(-2.5,1.8) (2.5,1.8)};
 \addplot[loss curve,domain=0:360,samples=100] ({2.40*cos(x)},{2.40*sin(x)});
\end{groupplot}
\end{tikzpicture}
\caption{The square map sends constant-product hyperbolas to torque lines and
loss circles to circles with squared radius. It does not remove the EESM's
third, torque-neutral coordinate.}
\label{fig:square-map}
\end{figure}


The price appears when voltage is restored. A general quadratic in
\((z_\psi,z_q,z_0)\) contains terms linear in products such as
\(z_\psi z_0\) and \(z_qz_0\). The square map represents
\(z_\psi^2-z_q^2\) and \(2z_\psi z_q\) elegantly, but not the individual
signed factors coupled to \(z_0\). Expressing them through
\cref{eq:square-root-inverse} introduces square roots and branches. Thus the
map linearizes torque superbly, preserves the loss radius, and provides a
memorable picture; it does not generally simplify the full EESM voltage
constraint enough to beat projective or performance-space coordinates. It is
retained as an interpretation and regression route, not selected as the
online solver.

\subsection{Why stereographic projection is also not selected}

After whitening, current directions lie on a unit sphere. Stereographic
projection sends a point \((x_1,x_2,x_3)\) away from the north pole to
\begin{equation}
 (s,t)=\left(\frac{x_1}{1-x_3},\frac{x_2}{1-x_3}\right),
 \quad
 \vect x=\frac{1}{1+s^2+t^2}
 \begin{bmatrix}2s\\2t\\s^2+t^2-1\end{bmatrix}.
\end{equation}
It is conformal and maps circles on the sphere to circles or lines. However,
substitution into either quadratic performance coordinate produces a quartic
numerator over \((1+s^2+t^2)^2\). Clearing the denominator gives quartics and
does not simultaneously diagonalize torque and voltage. The method is useful
for plotting every direction on one plane, but the excluded projection pole
requires a second chart and the algebra is no smaller. This elegant route is
therefore a documented dead end for the embedded solve.

\paragraph{One structure, three lenses.}
Projective \(\eta\) is the factor ratio. Hyperbolic \(\alpha\) is half its
logarithm. The square map stores the factor product in one Cartesian
coordinate and their difference of squares in the other. None is a mysterious
trick; all three describe the same constant-product geometry.
